\section*{Capítulo \ref{ch:b1:organic-redox} --- Oxidación y reducción en química orgánica}

\begin{solution}{exo:b1:organic-redox:1}
Etanol: \ce{CH3} $-3$, \ce{CH2OH} $-1$. Etanal: \ce{CH3} $-3$, CHO $+1$.
Ácido etanoico: \ce{CH3} $-3$, COOH $+3$. Propanona: \ce{CH3} $-3$ (dos
veces), C=O $+2$. Ácido metanoico: $+2$.
\end{solution}

\begin{solution}{exo:b1:organic-redox:2}
Propan-1-ol: propanal y, después, ácido propanoico. Propan-2-ol:
propanona. 2-Metilpropan-2-ol: no reacciona (el color naranja persiste).
\end{solution}

\begin{solution}{exo:b1:organic-redox:3}
\ce{CH3COCH3 + 2H+ + 2e- -> CH3CH(OH)CH3};
\ce{CH3CHO + 2H+ + 2e- -> CH3CH2OH}.
\end{solution}

\begin{solution}{exo:b1:organic-redox:4}
\ce{NaBH4}: butan-1-ol, butan-2-ol, no reacciona, no reacciona (el ácido
solo se desprotona). \ce{LiAlH4}: butan-1-ol, butan-2-ol, butan-1-ol y
etanol, butan-1-ol.
\end{solution}

\begin{solution}{exo:b1:organic-redox:5}
\ce{5CH3CH(OH)CH3 + 2MnO4- + 6H+ -> 5CH3COCH3 + 2Mn^2+ + 8H2O}.
$n = 1.20/60.0 = \qty{0.0200}{mol}$; permanganato, $\frac25 \times 0.0200 =
\qty{8.0e-3}{mol}$, \qty{400}{mL} de la disolución de \qty{0.020}{mol/L}.
\end{solution}

\begin{solution}{exo:b1:organic-redox:6}
Se calienta el alcohol con el \omterm{def:b1:nernst:couple}{oxidante} en un matraz montado para
destilar, a una temperatura entre las dos temperaturas de ebullición: el
propanal (\qty{48}{\celsius}) destila a medida que se forma y escapa a la
oxidación posterior; el propan-1-ol (\qty{97}{\celsius}) se queda en el
matraz.
\end{solution}

\begin{solution}{exo:b1:organic-redox:7}
Un par B--H del \ce{BH4-} ataca el carbono carbonílico, y el par $\pi$ del
C=O pasa al oxígeno; el \omterm{def:b1:alcohol-activation:alkoxide}{alcóxido} toma un protón del etanol. El hidrógeno
del carbono del producto (\ce{CH3CH2OH}, uno de los dos del C1) procede
del reactivo; el del oxígeno, del disolvente.
\end{solution}

\begin{solution}{exo:b1:organic-redox:8}
El butan-2-ol es \omterm{def:b1:stereochemistry-in-depth:stereocentre}{quiral}, pero la cetona plana es atacada por igual por
sus dos caras: una \omterm{def:b1:stereochemistry-in-depth:enantiomers}{mezcla racémica}, ópticamente inactiva.
\end{solution}

\begin{solution}{exo:b1:organic-redox:9}
\ce{3CH3CH2OH + 2Cr2O7^2- + 16H+ -> 3CH3COOH + 4Cr^3+ + 11H2O}: el
dicromato naranja (cromo $+$VI) se convierte en cromo(III) verde.
\end{solution}

\begin{solution}{exo:b1:organic-redox:10}
2-Metilbutan-2-ol: el etanal + \ce{CH3CH2MgBr} da butan-2-ol (una
adición, sin redox); oxidación a butanona; \ce{CH3MgBr} sobre la
butanona y, después, tratamiento. Una oxidación. Pentan-3-ol: el
propanal + \ce{CH3CH2MgBr} en una sola etapa, sin redox.
\end{solution}

\begin{solution}{exo:b1:organic-redox:11}
1. Se protege la cetona como \omterm{def:b1:acetals-protection:hemiacetal}{acetal} cíclico (etano-1,2-diol, ácido,
Dean--Stark). 2. El \ce{LiAlH4} en éter seco reduce el éster al alcohol
primario (y metanol). 3. El ácido acuoso elimina el \omterm{def:b1:acetals-protection:hemiacetal}{acetal}:
5-hidroxipentan-2-ona, \ce{CH3CO(CH2)3OH}.
\end{solution}

\begin{solution}{exo:b1:organic-redox:12}
\ce{RCH2OH}: carbono carbinólico a $-1$; \ce{RCOOH}: $+3$; cuatro
electrones, como en \ce{RCH2OH + H2O -> RCOOH + 4H+ + 4e-}. Del metano
($-4$) al dióxido de carbono ($+4$): ocho electrones,
\ce{CH4 + 2H2O -> CO2 + 8H+ + 8e-}.
\end{solution}

\begin{solution}{pb:b1:organic-redox:1}
\textbf{1.} C1 (el que lleva el OH): 0; los cinco \ce{CH2}: $-2$.
\textbf{2.} C1 (C=O): $+2$; los cinco \ce{CH2}: $-2$.
\textbf{3.} Los dos carbonos COOH: $+3$; los cuatro \ce{CH2}: $-2$.
\textbf{4.} El C1 ($0 \to +3$) y su vecino C6 ($-2 \to +3$), que se
convierte en el segundo carbono carboxílico; se rompe el enlace C1--C6.
\textbf{5.} $3 + 5 = 8$ electrones.
\textbf{6.} El ácido nítrico es lo bastante fuerte para romper el enlace
C--C contiguo al carbonilo, cosa que no hacen los \omterm{def:b1:nernst:couple}{oxidantes} más suaves.
\textbf{7.} \ce{C6H12O + 3H2O -> C6H10O4 + 8H+ + 8e-}.
\textbf{8.} $+$V y $+$I.
\textbf{9.} \ce{2HNO3 + 8H+ + 8e- -> N2O + 5H2O}.
\textbf{10.} Al sumar, se cancelan ocho electrones, \ce{8H+} y tres
moléculas de agua: \ce{C6H12O + 2HNO3 -> C6H10O4 + N2O + 2H2O}.
\textbf{11.} 8.
\textbf{12.} $10^6/146.0 = \qty{6.85e3}{mol}$ de ácido; $2 \times 6.85 \times
10^3 \times 63.0 = \qty{8.6e5}{g}$, \qty{0.86}{t} de ácido nítrico.
\textbf{13.} \ce{H2O2 + 2H+ + 2e- -> 2H2O}.
\textbf{14.} \ce{C6H12O + 4H2O2 -> C6H10O4 + 5H2O}.
\textbf{15.} $4 \times 6.85 \times 10^3 \times 34.0 = \qty{9.3e5}{g}$,
\qty{0.93}{t}.
\textbf{16.} El agua: ningún óxido de nitrógeno.
\textbf{17.} Una constante de equilibrio grande no dice nada de la
velocidad; el \ce{H2O2} reacciona lentamente sin \omterm{def:b1:elementary-steps:catalyst}{catalizador}, como ocurre
con su propia descomposición.
\textbf{18.} Reduce la ciclohexanona de nuevo a ciclohexanol:
\ce{4C6H10O + BH4- -> B(OC6H11)4-} y, después,
\ce{B(OC6H11)4- + 4H2O -> 4C6H11OH + B(OH)4-}.
\textbf{19.} \qty{6.85e3}{mol}.
\textbf{20.} \qty{6.85e3}{mol} de \ce{N2O} (uno por cada ácido).
\textbf{21.} $6.85 \times 10^3/0.95 \times 100.0 = \qty{7.2e5}{g}$, \qty{0.72}{t}.
\textbf{22.} $6.85 \times 10^3 \times 44.0 = \qty{3.0e5}{g}$: \textbf{\qty{0.30}{t}
de \ce{N2O} por tonelada de ácido adípico}.
\end{solution}
